% Standalone report fragment. Requires amsmath/amssymb; no new theorem
% environments or macros are declared. MP196 implements the unshifted
% target. The positive-target construction is a command-level alternative.
\section{Moving contradiction detection into modus ponens}
\label{sec:mp-halt}

The following command transformation keeps the numerical proof calculus
unchanged. At command level it replaces Shift's halting test with a test
made only after a successful modus-ponens comparison. There are two versions.
The 196-state construction uses the original target $d+x$ and reuses the
existing destructive comparator for the terminal singleton test. The
positive-target alternative uses $d+x+1$, so that the terminal test can
be a zero test of the retained offset $d+x$. Its command-level proof is
included as a separate possible refinement. Neither construction here
is claimed to improve the supplied 194-state bound.

This is a statement about refutation power and command compilation. The
number of binary-machine controls is a separate question, settled only by
a complete literal table and its source simulation. In particular, this
section does not infer a state saving from a shorter command description.

\subsection{Numerical stores and the retained axiom command}

Write
\[
 p(a,b)=\frac{(a+b)(a+b+1)}2+b,
 \qquad [a;w]=p(a,w),\qquad
 I(a,b)=4p(a,b),\quad U(i,a)=4p(a,i)+2.
\]
The total Cantor projections satisfy
\[
 \operatorname{hd}[a;w]=a,
 \quad\operatorname{tl}[a;w]=w,
 \quad[0;0]=0,
 \quad\operatorname{hd}(0)=\operatorname{tl}(0)=0.
\]
Thus stacks are natural numbers with total projections; they are not
assumed to be free lists. The constructor alias $I(0,0)=0$ is retained.
Code zero is the fixed theorem of truth, and code one is the designated
contradiction. A command store is $\langle t,w,d,h\rangle$, consisting of
the visible formula, stack, numerical offset, and cached formula.

The ten-entry separate-diagonal AX command is unchanged. With
$D=d+x$, $H=h+y$, its effect is
\[
 \mathsf{AX}_{A,x,y}\langle t,w,d,h\rangle
 =\langle g(H,D,A),[q_1,\ldots,q_9,t;w],0,0\rangle,
 \tag{MH1}
\]
where
\[
 (q_1,\ldots,q_9)=
 (\Gamma(A),L_6,L_5,L_4,L_3,L_2,L_1,
       \rho(D,A,H,d),U(d,h)).
\]
Here $L_1,\ldots,L_6$, $\Gamma$, and the guarded separation gate $g$
are the retained numerical axiom families, and
\[
 \rho(D,A,H,d)=
 I\bigl(I(I(H,d),D),I(I(D,H),I(A,H))\bigr),\qquad d\le D.
\]
Each $q_1,\ldots,q_8$ and $g$ is divisible by four, including a failed
guard's output zero. The ninth head satisfies $q_9\equiv2\pmod4$ and
$q_9\ge2$. None of these generated entries is one. The saved old top
$t$ is the tenth entry and is not subject to that residue assertion.

\subsection{The two new modus-ponens protocols}

In both protocols, Shift is the total command
\[
 \mathsf S\langle t,w,d,h\rangle
   =\langle h,\operatorname{tl}(w),d,\operatorname{hd}(w)\rangle.
 \tag{MH2}
\]
Equation (MH2) specifies the abstract protocol used for the compiler
identities below. MP196 agrees with it on every store reached from
blank-tape initialization. Its behavior on other stores retains a guard,
documented after the soundness proof; the distinction
is necessary when stating a literal simulation contract.
For the original-target command implemented by MP196, put $b=d+x$ and
define
\[
 \mathsf{MP}^{0}_{x}\langle t,w,d,h\rangle=
 \begin{cases}
  \mathsf{halt},&I(t,b)=h\text{ and }b=1,\\
  \langle b,w,0,0\rangle,&I(t,b)=h\text{ and }b\ne1,\\
  \langle0,w,b,0\rangle,&I(t,b)\ne h.
 \end{cases}
 \tag{MH3}
\]
Successful target zero returns the zero theorem; it does not halt.
The MP196 implementation reuses two existing source routines. At entry
to the successful tail, $T=H=X=0$ and $D=b$. The existing transfer loop
\texttt{node51} changes this to $T=b$, $X=b$, $D=H=0$. Increment $H$
once and invoke the existing \texttt{node7}, which destructively compares
$H$ with $X$ and halts on equality. The comparison now succeeds exactly
when $b=1$. Otherwise it clears $H,X$ and returns with $T=b$, $D=0$,
including when $b=0$. The exact new structured tail is
\[
 \bigl(\texttt{node51},
       (\mathsf{inc}\,H,\texttt{node7})\bigr).
\]
Only its two new binary readers are needed; the transfer and halting
comparator are retained. The checked source exit returns to the command
loop, so it realizes (MH3), not merely its register arithmetic.

For the positive-target version, put $k=d+x$ and $b=k+1$, and let
$r\colon\mathbb N\to\mathbb N$ be any fixed total output-offset function:
\[
 \mathsf{MP}^{+}_{x}\langle t,w,d,h\rangle=
 \begin{cases}
  \mathsf{halt},&I(t,b)=h\text{ and }k=0,\\
  \langle b,w,r(k),0\rangle,&I(t,b)=h\text{ and }k>0,\\
  \langle0,w,k,0\rangle,&I(t,b)\ne h.
 \end{cases}
 \tag{MH4}
\]
Examples include $r(k)=k$, $r(k)=0$, and $r(k)=k\mathbin{\dot-}1$,
where $\dot-$ is truncated subtraction. The proof compiler immediately clears the successful
offset using a total Pop word, so every such convention has the same
proof. The target $b=1$ is recognized by one zero test of $k$,
provided $I(t,k+1)$ is formed while retaining $k$. The algebraic compiler
identities use only that $r(k)$ is a natural number. An arithmetized
machine implementation must additionally supply a terminating source
routine for its chosen $r$ with PA-provable correctness, as is immediate
for the three listed elementary functions. The delivered MP196 table
implements (MH3), not this positive-target alternative.

It is essential that the singleton test occurs only after the equality
comparison succeeds. A failed comparison with proposed target one does
not establish a contradiction and must return normally.

\subsection{Uniform stack identities}

Let $A_0=\mathsf{AX}_{0,0,0}$, and let $\mathsf{MP}^{\epsilon}$ stand
for either of the preceding protocols. Define
\[
 \begin{aligned}
 \mathsf{Push}&=A_0;\mathsf S^9;\mathsf{MP}^{\epsilon}_0,\\
 \mathsf{Pop}&=A_0;\mathsf S^{11},\\
 \mathsf{Discard}&=\mathsf{MP}^{\epsilon}_0,\\
 \mathsf{Offset}_i&=\mathsf{MP}^{\epsilon}_i.
 \end{aligned}
 \tag{MH5}
\]
These words have the exact contracts
\[
 \begin{aligned}
 \mathsf{Push}\langle t,w,d,h\rangle
    &=\langle0,[t;w],0,0\rangle,\\
 \mathsf{Pop}\langle t,[a;w],d,h\rangle
    &=\langle t,w,0,a\rangle,\\
 \mathsf{Discard}\langle t,w,0,0\rangle
    &=\langle0,w,0,0\rangle,\\
 \mathsf{Offset}_i\langle0,w,0,0\rangle
    &=\langle0,w,i,0\rangle.
 \end{aligned}
 \tag{MH6}
\]
Push and Pop quantify over all incoming stores, including arbitrary
cached values. Discard's displayed clean-cache precondition is sufficient
for the compiler. For (MH3) it can be strengthened to arbitrary $h$,
since the success and failure branches at target zero both return
$\langle0,w,0,0\rangle$. That stronger assertion is not made for (MH4).

\paragraph{Proof of the identities.}
After $A_0;\mathsf S^9$ the store is
\[
    \langle q_8,[t;w],0,U(d,h)\rangle.
\]
The final MP comparison is between an implication of residue zero and
a universal formula of residue two, so it fails. Its resulting offset
is zero in either protocol. This proves Push. After ten Shifts following
$A_0$, the store is $\langle U(d,h),w,0,t\rangle$; one more unconditional
Shift proves Pop. This includes
$\mathsf{Pop}\langle t,0,d,h\rangle=\langle t,0,0,0\rangle$.

In (MH3), Discard proposes target zero, and both outcomes coincide.
In (MH4), it compares $I(t,1)>0$ to cache zero and therefore fails.
For Offset, the positive-target version compares $I(0,i+1)>0$ to zero.
The unshifted version fails when $i>0$ and returns zero in either branch
when $i=0$. The resulting offset is exactly $i$. All these calculations
respect $[0;0]=0$ and $I(0,0)=0$. No distinction between a zero push onto
zero and an empty canonical stack is used. $\square$

\subsection{Selectors and generalization}

The unchanged AX stack order gives the same selectors as in the no-push
PA194N compiler. Every selector below starts at a clean-cache store
$\langle t,w,0,0\rangle$ and preserves its arbitrary underlying stack.
For generated head $q_j$, $1\le j\le7$, use
\[
 \mathsf{Discard};\mathsf{AX}_{A,D,H};
       \mathsf S^{j+1};\mathsf{Pop}^{9-j}.
 \tag{MH7}
\]
After the initial Shifts, the unselected generated tail and the saved
zero remain. The Pops preserve the selected top, remove that tail, and
finally put the saved zero into the cache. The output is
$\langle q_j,w,0,0\rangle$.

For the dominated propositional head $\rho(a,b,c,u)$, $u\le a$, use
\[
 \mathsf{Discard};\mathsf{Offset}_u;
   \mathsf{AX}_{b,a-u,c};\mathsf S^9;\mathsf{Pop}.
 \tag{MH8}
\]
The old offset is $u$ and the new parameter is
$u+(a-u)=a$, so the required instance is selected literally. A legitimate
guarded gate is selected by
\[
 \mathsf{Discard};\mathsf{AX}_{A,D,H};\mathsf{Pop}^{10}.
 \tag{MH9}
\]
Truth zero uses Discard alone.

Generalization at numerical label $i$ uses
\[
 \mathsf{Gen}_i=
 \mathsf{Push};\mathsf{Offset}_i;\mathsf S;
      \mathsf{AX}_{i,0,0};\mathsf S^{10}.
 \tag{MH10}
\]
From $\langle a,w,0,0\rangle$, its state immediately before AX is
$\langle0,w,i,a\rangle$. The ninth generated head is $U(i,a)$ and the
saved top is zero, so the final ten Shifts give
$\langle U(i,a),w,0,0\rangle$. These words do not invoke a new axiom or
modify a numerical formula. They only select and manipulate the retained
axiom instances and theorem codes.

\subsection{A complete finite-proof compiler}

Compile an axiom leaf with its selector and a generalization node with
(MH10). Consider an MP proof node with major conclusion $I(a,b)$ and
minor conclusion $a$. For $b>0$, use
\[
 \begin{split}
 \mathsf{Discard};\mathsf{Push};C(\text{major});\mathsf{Push};
 C(\text{minor});\mathsf{Push};\mathsf S;\mathsf S;\qquad\qquad\\[-2pt]
 \mathsf{MP}^{\epsilon}_{b-\epsilon};\mathsf{Pop},
 \end{split}
 \tag{MH11}
\]
where $\epsilon=0$ for (MH3) and $\epsilon=1$ for (MH4).
If either recursive call has already halted, the whole word halts.
Otherwise, immediately before the two Shifts the store is
\[
       \langle0,[a,I(a,b),0;w],0,0\rangle.
\]
The two total Shifts produce
$\langle a,[0;w],0,I(a,b)\rangle$. The MP comparison therefore succeeds.
For $b=1$ it halts immediately. For $b>1$ it returns $b$ with stack
$[0;w]$ and zero cache. Its offset is zero in (MH3) and $r(b-1)$ in
(MH4), which is harmless because the Pop contract permits every offset.
The total Pop contract removes exactly the reserved zero and returns
$\langle b,w,0,0\rangle$.

If $b=0$, the compiler may use Discard without compiling either premise.
Code zero already has a fixed derivation, so this shortcut preserves the
conclusion. Thus the positive-target restriction removes no refutation
power. The unshifted protocol also allows direct execution of successful
MP with target zero; its contract is explicitly checked, but the compiler
does not need it.

Induction on a finite numerical derivation $e$ gives two statements:
\begin{enumerate}
\item From every $\langle t,w,0,0\rangle$, $C(e)$ either halts or returns
      $\langle\operatorname{concl}(e),w,0,0\rangle$.
\item If $\operatorname{concl}(e)=1$, the word $C(e)$ halts.
\end{enumerate}
For the second statement, an axiom leaf cannot have code one by the
residue facts in (MH1), and a generalization has residue two. An MP node
with conclusion one halts at its final successful comparison unless a
subproof has already halted. Hence no final Shift test is needed.

For every refutation $e$ and every natural $k$, the word
\[
        W_{e,k}=\mathsf{AX}_{k,0,0};C(e)
 \tag{MH12}
\]
halts from every incoming store. The first AX clears both caches, after
which the compiler's clean-entry theorem applies. This assertion is
deliberately stronger than compilation from an empty stack: the scheduler
retains its store between unsuccessful candidates.

\subsection{Explicit scheduling and the cofinal witness family}

The unchanged command decoder uses positive continuation pointers.
For finite words define
\[
 \begin{aligned}
 e([])&=0,\\
 e(\mathsf S::v)&=1+p(0,e(v)),\\
 e(\mathsf{MP}^{\epsilon}_x::v)&=1+p(1,p(x,e(v))),\\
 e(\mathsf{AX}_{c,x,y}::v)&=1+p(c+2,p(y,p(x,e(v)))).
 \end{aligned}
 \tag{MH13}
\]
At a fetch boundary $(n,l)$, the scheduler decodes payload $n$ and
increments $n$ when $l=0$; otherwise it decodes $l-1$ and keeps $n$.
The decoded residual pointer of any nonempty word in (MH13) is exactly
the code of its tail. For arbitrary positive $l$, the residual is at
most $l-1$, since each Cantor projection is at most its input. Thus each
arbitrary candidate reaches its next fresh boundary after finitely many
commands, or halts first. No reset of the command store is required.

The fresh candidates from the initialized scheduler are $2,3,4,\ldots$.
For (MH12), the exact fresh candidate is $e(W_{e,k})-1$, and
\[
             e(W_{e,k})-1\ge k+2.
 \tag{MH14}
\]
The witness family is therefore cofinal: from any current fresh counter
$n$, take $k\ge n$. Unless an earlier candidate has halted, bounded
induction through the intervening candidates reaches this witness, whose
all-store property then forces a halt. This supplies the scheduling
argument explicitly, instead of appealing to an unspecified fairness
assumption. The proof is arithmetizable in PA using induction on positive
continuation codes and the finite interval of fresh candidate numbers.

\subsection{Soundness from the initialized store}

The soundness statement has a different quantifier from (MH12).
It starts with the initialized theorem-annotated store, not with an
arbitrary store containing possibly unproved formulas. Annotate the
visible top, cached formula, and every live canonical stack entry with
derivations in the retained numerical calculus. The initialized zero
entries carry the fixed zero derivation. Popping numerical zero uses that
same derivation; pushing zero onto zero creates no extra canonical entry.

AX preserves the annotations: its heads and gate are axiom instances,
its saved top retains its old derivation, and its ninth head is a
generalization of the cached theorem. Unconditional Shift only moves
annotated formulas. Failed MP returns the zero theorem. Successful
nonhalting MP returns its conclusion with the usual MP derivation.
The sole halting event is a successful MP comparison whose conclusion
is one. It therefore extracts a numerical-calculus refutation.

Together with the compiler and scheduler argument, this proves
equivalence of halting of the specified initialized command scheduler
and refutability in the unchanged numerical calculus. The proof transforms
finite syntactic objects; it does not invoke a truth predicate or
reflection principle. The constructors, stack operations, proof compiler,
and proof annotations are primitive recursive, and PA proves their
correctness by the displayed finite inductions. Composition with the
inherited ordinary PA calibration gives the corresponding source-level
consistency equivalence.

\subsection{The retained Shift guard and its invariant}

MP196 retains the original Shift source block: on a store with
$t+h\ne1$ it performs (MH2), and on a store with $t+h=1$ it halts.
The retained halting branch is unreachable on the new initialized run,
as shown below. The literal table can therefore reuse the same
\texttt{node7} comparator in the successful MP tail while retaining
complete source binding, including the unused old Shift halt leaf.
Its exact all-store contract is
\[
 \mathsf S_{\mathrm{literal}}\langle t,w,d,h\rangle=
 \begin{cases}
  \mathsf{halt},&t+h=1,\\
  \langle h,\operatorname{tl}(w),d,\operatorname{hd}(w)\rangle,&t+h\ne1.
 \end{cases}
 \tag{MH15}
\]
Consequently, an unconditional all-store simulation claim for (MH2)
would be false for the literal table.

The needed restriction is an elementary numerical invariant, independent
of consistency or of whether any formula has a proof. Call a command
store admissible when its top, cache, and every live canonical stack
entry differ from one. The initialized store is admissible. Every newly
generated AX entry differs from one by its constructor residue; its saved
top was already admissible. Failed MP returns zero. Successful MP in
either version halts with target one, or returns a target different
from one. An ordinary Shift permutes entries of the admissible store,
with a zero returned by an exhausted canonical stack. Finally, for
natural numbers $t,h\ne1$, the equality $t+h=1$ is impossible.

Induction over completed commands therefore proves that every prehalting
command boundary is admissible and that the old Shift halting branch
is never used. In particular, this is not an assumption that PA is
consistent: a refuting run halts inside its successful MP before returning
a code one at a completed command boundary. The physical MP tail does
temporarily store one while performing that final test; the invariant is
asserted at completed command boundaries, not at every internal source
instruction. The literal source simulates the total-Shift protocol on its
entire initialized run.

The all-store compiler theorem for the abstract protocol remains valid
as stated. For the physical completeness direction, it is sufficient to
apply it at the admissible store actually reached by the scheduler when
the cofinal witness begins. All intermediate nonhalting stores remain
admissible, so the literal and abstract runs continue to agree until the
compiled successful MP halts. Arbitrary incoming stores are used only
in the abstract compiler identities; they are not silently included in
the literal Shift simulation contract.

To transfer that result to a literal binary table, one still needs the
exact initialization, complete source binding, and positive finite
simulation for every primitive source outcome. Those are machine-specific
obligations. Once supplied, a source halt gives a physical halt by
concatenating simulation segments. Conversely, a first physical halt at
time $T$ contradicts $T+1$ nonhalting source steps and their positive
simulation segments, by determinism. This is the usual PA-internal
finite-trace argument; the command proof alone does not supply its table
equalities or its state count.

\subsection{Executable checks and proof boundary}

The accompanying checker verifies symbolic identities over free store,
stack, and parameter terms. Its comparison evaluator accepts only
syntactic equality, explicit zero identities, provable positivity, and
incompatible constructor residues. Unresolved symbolic comparisons raise
an error; they are not guessed. The unshifted target-zero case uses an
explicit join of its two identical outcomes. Separate integer checks
execute Cantor pairing and total inversion, and the continuation-code
checks use a separate square-root decoder. The symbolic identities are
the uniform part of this audit; the integer examples are additional
implementation checks.

This is an ordinary mathematical proof with executable algebraic checks,
not a newly compiled end-to-end Lean or Rocq theorem. The inherited PA
calibration and the literal-machine simulation remain explicitly separate
dependencies. In particular, neither a successful command audit nor an
optimized source word is described as a certified sub-194-state machine.
